% Regression: dir= rides a resolved physical index.  An internal directed
% contraction and directed open ends on all four bearings keep the typed
% record, the physical-leg stroke, and the direction mark; the reversed
% spelling keeps the reversed barb.
\documentclass{standalone}
\usepackage{tenkz}
\makeatletter
\ExplSyntaxOn
\file_input:n { tenkz-kernel.code.tex }
\int_new:N \g__tenkz_test_dir_physical_stroke_int
\int_new:N \g__tenkz_test_dir_marks_int
\int_new:N \g__tenkz_test_dir_marks_reversed_int
\cs_new_eq:NN
  \__tenkz_test_wire_stroke:nn
  \__tenkz_kernel_r_wire_stroke:nn
\cs_set_protected:Npn \__tenkz_kernel_r_wire_stroke:nn #1#2
  {
    \exp_args:NV \__tenkz_model_get:nnN
      \l__tenkz_kernel_r_wire_record_tl {dir} \l_tmpa_tl
    \quark_if_no_value:NF \l_tmpa_tl
      {
        \str_if_eq:eeF {#1} {physical~leg}
          {\errmessage{directed~physical~index~lost~physical~stroke}}
        \int_gincr:N \g__tenkz_test_dir_physical_stroke_int
        \str_case:Vn \l_tmpa_tl
          {
            {to}
              {
                \tl_if_in:NnF \l__tenkz_kernel_r_wire_ink_tl {dir~marks}
                  {\errmessage{directed~physical~index~lost~its~barb}}
                \int_gincr:N \g__tenkz_test_dir_marks_int
              }
            {from}
              {
                \tl_if_in:NnF \l__tenkz_kernel_r_wire_ink_tl
                  {dir~marks~reversed}
                  {\errmessage{reversed~physical~index~lost~its~barb}}
                \int_gincr:N \g__tenkz_test_dir_marks_reversed_int
              }
          }
      }
    \__tenkz_test_wire_stroke:nn {#1} {#2}
  }
\AtEndDocument
  {
    \int_compare:nNnF { \g__tenkz_test_dir_physical_stroke_int } = {6}
      {\errmessage{directed~physical~stroke~count~changed}}
    \int_compare:nNnF { \g__tenkz_test_dir_marks_int } = {5}
      {\errmessage{directed~physical~barb~count~changed}}
    \int_compare:nNnF { \g__tenkz_test_dir_marks_reversed_int } = {1}
      {\errmessage{reversed~physical~barb~count~changed}}
  }
\ExplSyntaxOff
\makeatother
\begin{document}
\tenkzkernel
% Open directed physical ends on all four bearings around one site.
\begin{tenkz}[rows={wire,wire,wire}, cols=3, bonds=none]
  \tn[at=(2,2), name=X, skin=dot,
    ports={0:physical, 90:physical, 180:physical, 270:physical}]{}
  \tnwire[dir=to, name=east]{X.0}{open e}
  \tnwire[dir=to, name=north]{X.90}{open n}
  \tnwire[dir=to, name=west]{X.180}{open w}
  \tnwire[dir=to, name=south]{X.270}{open s}
\end{tenkz}
% An internal directed physical contraction: two consumed ports, no
% boundary entry.
\begin{tenkz}[rows={wire}, cols=2, bonds=none]
  \tn[at=(1,1), name=A, ports={0:physical}]{A}
  \tn[at=(1,2), name=B, ports={180:physical}]{B}
  \tnwire[dir=to, name=inner]{A.0}{B.180}
\end{tenkz}
% The reversed spelling on an open end: the barb flips, the index enters.
\begin{tenkz}[rows={wire}, cols=1, bonds=none]
  \tn[at=(1,1), name=C, ports={90:physical}]{C}
  \tnwire[dir=from, name=entering]{C.90}{open n}
\end{tenkz}
\end{document}
